\section{Implementación}

\subsection*{Estructura}

Dado el alcance del proyecto, la organización del repositorio se realizó mediante un \textbf{monorepo}. Lo anterior,
permite centralizar el código tanto del cliente, como del servidor, y orquestar tareas comúnes de desarrollo.

\dirtree{%
    .1 ./.
    .2 client/.
    .3 cmd/.
    .4 client/.
    .3 internal/.
    .4 domain/.
    .4 network/.
    .4 protocol/.
    .4 ui/.
    .5 messages/.
    .5 panels/.
    .5 styles/.
    .2 docs/.
    .3 mockups/.
    .3 report/.
    .4 figures/.
    .4 sections/.
    .3 tapes/.
    .2 server/.
    .3 src/.
    .4 protocol/.
    .2 stress/.
}

Por un lado, el \textbf{Client} se organiza siguiendo la
\href{https://github.com/golang-standards/project-layout/blob/master/README_es.md}{Standard Go
Project Layout}, la organización estándar de proyectos en \textbf{Go}. Utiliza \textbf{cmd} para
el punto de entrada del cliente e \textbf{internal} para organizar el código interno, relacionado
tanto con la lógica de negocio como con la vista de usuario.

Por otro lado, el \textbf{servidor} se organiza de una manera más simple, en un único
\textbf{crate} que centraliza el comportamiento general de la aplicación. Únicamente el protocolo
se separa en su propio módulo.

\subsection*{Dependencias}

Las dependencias del desarrollo se listan en la la tabla \ref*{tab:requisitos}.

\begin{table}[h]
    \centering
    \begin{tabular}{ll}
        \textbf{Componente} & \textbf{Versión} \\ \hline
        \textbf{Go} & 1.26.6 \\
        \textbf{Bubble Tea} & 1.3.10 \\
        \textbf{Lipgloss} & 1.1.0 \\
        \textbf{Bubbles} & 1.0.0 \\
        \textbf{Rust} & 1.97.1 \\
        \textbf{Tokio} & 1.53.1 \\
        \textbf{Serde} & 1.0.229 \\
        \textbf{Serde JSON} & 1.0.151 \\
        \textbf{Just} & 1.57.0 \\
        \textbf{Prek} & 0.5.0 \\
        \textbf{k6} & 2.3.0 \\
        \textbf{xk6-tcp} & main \\
    \end{tabular}
    \caption{Requisitos y dependencias del proyecto.}
    \label{tab:requisitos}
\end{table}

\subsection*{Flujo de Trabajo}

Para unificar los sistemas de compilación, ejecución de pruebas unitarias, documentación y ejecución de ambos
programas, se utilizó \textbf{Just}. Por ejemplo, el siguiente comando realiza el build tanto del cliente como del servidor.

\begin{lstlisting}[language=bash]
    just build
\end{lstlisting}

Además, para mantener una consistencia de código se utiliza \textbf{Prek} para realizar un formateo del código,
tanto de \textbf{Rust} como de \textbf{Go}, antes de cada commit.

La mayoría del flujo inicial se basó en el flujo de trabajo \textit{code a little, test a little}. Las pruebas unitarias
se ejecutan en \textbf{workflows} de \textbf{integración continua} durante cada push.

\subsection*{Servidor}

Observemos como se translada el modelo del servidor al código. 

\subsubsection*{Multiplexado}

La macro \texttt{tokio::select!} permite ejecutar concurrentemente varias expresiones en la misma \textbf{task}. De está manera, 
ambas fuentes de eventos alternan su ejecución dependiendo del cumplimiento de las ramas.

\begin{lstlisting}[language=Rust]
tokio::select! {
    // 1. Lectura del socket TCP
    result = reader.read_line(&mut line) => {
        match result {
            Ok(0) => break, // EOF
            Ok(_) => {
                if let Ok(msg) = serializer::deserialize(&line) {
                    if let Err(err) = route_msg(msg, &_username, &hub, &mut writer).await {
                        send_msg(&mut writer, new_response(TypeC2S::Invalid, err, None)).await;
                        break;
                    }
                }
            }
            Err(_) => break,
        }
    }

    // 2. Recepcion de mensajes salientes desde el canal MPSC
    Some(msg) = rx.recv() => {
        send_msg(&mut writer, msg).await;
    }

    else => break,
}
\end{lstlisting}

\subsubsection*{Enrutamiento}

La función de enrutamiento implementa el patrón de descrito en la \autoref{sec:analisis-diseno}: recibe un 
\texttt{ClientMessage} ya deserializado y, mediante un \texttt{match} exhaustivo sobre sus variantes, delega 
la operación correspondiente a una función \textit{handler} específica.

\begin{lstlisting}[language=Rust]
async fn route_msg<W>(
    msg: ClientMessage,
    username: &str,
    hub: &Arc<Mutex<Hub>>,
    writer: &mut W,
) -> Result<(), MessageResult>
where
    W: AsyncWrite + Unpin,
{
    match msg {
        ClientMessage::Users => handle_users_list(hub, writer).await,
        ClientMessage::Status { status } => handle_status(username, status, hub).await,
        ClientMessage::PublicText { text } => handle_public_text(username, &text, hub).await,
        ClientMessage::NewRoom { roomname } => {
            handle_create_room(username, &roomname, hub, writer).await?;
        }
        ...
    }
}
\end{lstlisting}

\subsubsection*{Estado Compartido}

La mayoría de funciones \texttt{handler} recibe un \texttt{\&Arc$<$Mutex$<$Hub$>$$>$}. Está es una manera popular en la comunidad
de \textbf{Rust} para implementar el \textbf{estado compartido por exclusión mutua} para realizar modificaciones en un estado global.
Sin embargo, la cautela al abrir el bloqueo tiene que realizarse minuciosamente y abarcando la menor cantidad de 
operaciones.

Por ejemplo, apoyanse ampliamente de \textit{scopes} resulta indispensable. 

\begin{lstlisting}[language=Rust]
let result = {
    let hub = hub.lock().unwrap();
    hub.send_to(&msg, sender, receiver)
    // El hub se libera al salir de este scope
};

match result {
    Ok(_) => {}
    Err(e) => {
        send_msg(
            writer,
            new_response(TypeC2S::Text, e, Some(receiver.to_string())),
        )
        .await;
    }
}
\end{lstlisting}

\subsection*{Cliente}

\subsubsection*{Concurrencia en el Cliente}

Mantener sincronizada la interfaz de usuario con el estado global del servidor sin bloquear el hilo de
renderizado es una tarea indispensable para aplicaciones con interfaces visuales. Por está razón, la 
concurrencia resulta igual de necesaria en la capa del \textbf{cliente} que en \textbf{servidor}; no quieres
ofrecer una experiencia de usuario lenta y poco responsiva.

La conexión se encuentra desacoplada de la interfaz y corre en una \textbf{goroutine} secundaria. Por cada
línea que lee del buffer, si no está malformada, se deserializa y se envía a un canal.

\begin{lstlisting}[language=Go]
// Listen escucha los mensajes del servidor al cliente y los manda a traves de un canal
func (manager *ConnectionManager) Listen() {
	for {
		line, err := manager.reader.ReadString('\n')
		if err != nil {
			log.Printf("Connection closed: %v", err)
			return
		}

		msg, err := protocol.Deserialize(line)

		if err != nil {
			log.Printf("Invalid message: %s", line)
			continue
		}

		manager.incoming <- msg
	}
}
\end{lstlisting}

El modelo principal espera mensajes de ese canal. Cuando recibe alguno, retorna un comando con el mensaje
ya deserializado.

\begin{lstlisting}[language=Go]
func (m Model) waitForServerMsg() tea.Cmd {
    return func() tea.Msg {
        return <-m.conn.Messages()
    }
}
\end{lstlisting}

El mensaje es capturado por \textbf{Update} y delegado al \textit{router} de mensajes del servidor. Una vez
procesado, el modelo se resubscribe al canal en espera de más mensajes.

\begin{lstlisting}[language=Go]
func (m Model) Update(msg tea.Msg) (tea.Model, tea.Cmd) {
	switch msg := msg.(type) {

	case protocol.ClientMessage:
		return m, m.routeClientMessage(msg)

	case protocol.ServerMessage:
		cmd := m.routeServerMessage(msg)
        // Se resuscribe al canal
		return m, tea.Batch(cmd, m.waitForServerMsg())

	return m, nil
}
\end{lstlisting}

\subsection*{Modales}

Los modales utilizados en la interfaz por terminal implementan la interfaz descrita en la
\autoref{fig:modals-uml}. Cada tipo de modal (\texttt{Input}, \texttt{List}, \texttt{Confirm},
\texttt{Notification}) satisface esta interfaz, lo que permite al modelo raíz tratarlos de manera
uniforme.

El modelo raíz contiene un campo \texttt{activeModal} que implementa la máquina de estados descrita
en el diseño: mientras este campo esté activo, se le delega la totalidad del flujo de entrada hasta
resolverse.

\begin{lstlisting}[language=Go]
type Modal interface {
    InputCapturer
    Init() tea.Cmd
    Update(msg tea.Msg) (Modal, tea.Cmd)
    View() string
    Focus() tea.Cmd
    SetSize(width, height int)
}
\end{lstlisting}

Por ejemplo, un evento puede invocar la creación de un modal y establecerlo en la vista principal.

\begin{lstlisting}[language=Go]
case "c":
    return m.openModal(panels.NewInputModal("Create Room", "Roomname", 16, func(value string) tea.Msg {
        room, _ := protocol.NewRoomMessage(value)
        return room
    }, true))
\end{lstlisting}

La función \texttt{openModal} realiza la transición de estado \textit{sin modal} a \textit{modal
activo}, y espera a que el modal se resuelva en algún momento posterior.

\begin{lstlisting}[language=Go]
func (m *Model) openModal(modal panels.Modal) tea.Cmd {
    modal.SetSize(m.width, max(0, m.heigth-3))
    m.chatModel.Blur()
    cmd := modal.Focus()
    m.activeModal = modal
    // Inicializa el modal
    return tea.Batch(cmd, modal.Init())
}
\end{lstlisting}

Cuando el servidor responde exitosamente a la operación solicitada por el modal, el enrutador de mensajes 
del servidor resuelve el flujo: cierra el modal activo
y actualiza el estado local con la nueva información.

\begin{lstlisting}[language=Go]
if msg.Result == protocol.SUCCESS {
    m.closeModal()

    room := m.session.AddRoom(domain.RoomChannel, msg.Extra)
    m.roomsModel.AddItem(NewRoomItem(room))
}
\end{lstlisting}

Finalmente, \texttt{closeModal} realiza la transición inversa, devolviendo la interfaz al estado sin
captura.

\begin{lstlisting}[language=Go]
func (m *Model) closeModal() {
    m.activeModal = nil
    m.syncFocus()
}
\end{lstlisting}

\subsection*{Resultados}

Agregado a las pruebas unitarias, se crearon pruebas de éstres mediante una extensión de \textbf{K6} personalizada
para TCP, \textbf{xk6-tcp}. Las pruebas simulan un flujo concurrente de alta carga: 

\begin{enumerate}
    \item Conecta \textbf{3000} usuarios distribuidos en varias etapas.
    \item Cada usuario envía 5 mensajes públicos y espera 10 segundos.
    \item El usuario se desconecta.
\end{enumerate}

Los resultados obtenidos, se muestran en la siguiente tabla \ref{tab:estres}.

\begin{table}[htbp]
    \centering
    \small
    \begin{tabular}{lrr}
        \hline
        \textbf{Métrica} & \textbf{Total} & \textbf{Promedio} \\ \hline
        Pico de usuarios concurrentes & 3,000 & -- \\
        Conexiones TCP totales & 17,608 & 143.14 conn/s \\
        Errores de conexión & 1 & -- \\
        Escrituras TCP & 122,030 & 992.00 op/s \\
        Lecturas TCP & 2,582,224 & 20,991.38 op/s \\
        Datos enviados por clientes & 7.1 MB & 58 kB/s \\
        Datos distribuidos por servidor & 583 MB & 4.7 MB/s \\
        \hline
    \end{tabular}
    \caption{Resultados de la prueba de estrés TCP.}
    \label{tab:estres}
\end{table}